\section{Introduction}\label{sec:introduction}\label{tree:sec:introduction}

Random planar maps model random two-dimensional geometry. A
statistical-mechanical weight changes the law of the map and, in the
continuum description, the parameter of Liouville quantum gravity (LQG).
Contour encodings and mating of trees identify the corresponding quantum
surfaces. The metric problem asks whether shortest-path distances in the
maps converge to the intrinsic distances on those surfaces.

We consider two finite spherical ensembles: critical
Fortuin--Kasteleyn (FK) maps and maps decorated by a spanning tree.
The FK proof applies to every fixed \(q\in(0,4)\). The spanning-tree
endpoint has a separate proof using independent walk increments.
Both conclusions retain quantum area through a natural discrete measure.
Here \emph{subcritical} refers to the LQG range \(\gamma<2\), or equivalently
\(q<4\); the FK model itself is at its self-dual critical point.

\subsection{The ensembles and the limit}

A rooted planar map is a finite connected multigraph embedded in the
oriented sphere, considered up to orientation-preserving homeomorphism,
with one distinguished oriented edge. Loops and multiple edges are
allowed. Fix a positive integer \(n\).

For the FK ensemble, let \(\mathcal F_n\) consist of rooted pairs
\((M,A)\), where \(M\) has \(n\) edges and \(A\subseteq E(M)\).
The spanning subgraph \((V(M),A)\) includes isolated vertices. In the
dual map, let \(A^*\) contain the edges crossing \(E(M)\setminus A\).
Writing \(k\) for the number of connected components, set
\[
 \ell(M,A)=k(A)+k(A^*)-1=2k(A)+|A|-|V(M)|.
\]
This counts the interfaces between primal and dual FK components.
For fixed \(q\in(0,4)\), sample according to
\begin{equation}\label{eq:fk-law}
 \Pbb[(M_n,A_n)=(M,A)]
  =\frac{q^{\ell(M,A)/2}}{Z_n(q)},
 \qquad Z_n(q)=\sum_{(M,A)\in\mathcal F_n}q^{\ell(M,A)/2}.
\end{equation}
For the spanning-tree ensemble, sample \((M_n,T_n)\) uniformly from
rooted map--spanning-tree pairs with \(n\) edges. Equivalently, its map
marginal is weighted by the number of spanning trees.

In either ensemble, \(d_n\) denotes graph distance using every primal
map edge, each of length one. Equip the vertex set with the probability
measure
\begin{equation}\label{eq:degree-measure}
                  \mu_n(v)=\frac{\deg_{M_n}(v)}{2n},
\end{equation}
where a loop contributes twice to the degree. This is the law of the
vertex incident to a uniformly sampled corner of the map.

For FK maps, choose the unique \(\gamma\in(\sqrt2,2)\) satisfying
\begin{equation}\label{eq:q-gamma}
                 q=2+2\cos(\pi\gamma^2/2).
\end{equation}
For spanning-tree maps, set \(\gamma=\sqrt2\). The FK inventory parameter
and Brownian correlation are
\[
 b_{\mathrm F}=\frac{\sqrt q}{2+\sqrt q},\qquad
 \rho=\frac{b_{\mathrm F}}{1-b_{\mathrm F}}
      =\frac{\sqrt q}{2}=-\cos(\pi\gamma^2/4).
\]
We reserve \(p=1/d_\gamma\) for the reciprocal metric dimension; it is
not the inventory parameter. These conventions are used throughout.
Let \((S,h)\) be the
ordinary unit-area \(\gamma\)-quantum sphere: the quantum-sphere measure
of \cite[Definition~4.21(ii)]{DMS}, disintegrated at area one, with
its distinguished points forgotten. Write \(\mu_h\) for its quantum-area
probability measure and \(D_h\) for its intrinsic LQG metric, with one
fixed deterministic multiplicative normalization
\cite{GwynneMillerMetric,GwynneMillerConformal}.

\begin{theorem}[Metric and measure convergence]\label{thm:main}\label{tree:thm:main}
For each fixed \(q\in(0,4)\) in the FK ensemble, and for the spanning-tree
ensemble, there are deterministic positive numbers \(a_n\to0\) such that
\[
          (V(M_n),a_nd_n,\mu_n)
             \xrightarrow[n\to\infty]{\mathrm{law}}
          (S,D_h,\mu_h)
\]
in Gromov--Hausdorff--Prokhorov topology on compact metric probability
spaces modulo measure-preserving isometry. The limit has the parameter
specified above. Convergence holds through all positive integers \(n\);
the root and decoration are forgotten.
\end{theorem}

Gromov--Hausdorff--Prokhorov distance simultaneously compares the metric
spaces and their measures after isometric embedding in a common space;
we give the precise definition in Section~\ref{measure:section}.
In particular, the theorem implies convergence in the ordinary
Gromov--Hausdorff topology. The normalization may depend on the chosen
ensemble and on the fixed FK parameter. No uniformity in \(q\), endpoint
FK statement, normalization power law, or convergence rate is asserted.
The measure in Equation~\eqref{eq:degree-measure} is part of the theorem;
a conclusion for the uniform measure on vertices would require a
separate argument.

\subsection{History and significance}

The FK weights originate in the random-cluster model of Fortuin and
Kasteleyn \cite{FortuinKasteleyn1972}. Mullin's bijection
\cite{Mullin1967} encodes maps with a spanning tree by walks in a
quadrant; Bernardi gives a parenthesis-shuffle formulation
\cite{Bernardi2007}, and the precise finite convention used here appears
in \cite[Section~2.3, Theorem~1]{DHBijection}.
Sheffield's inventory bijection gives the FK encoding and its bilateral
Brownian limit \cite{SheffieldInventory}; Gwynne and Sun established
the finite empty-word limit \cite{GwynneSunFinite}.
The quantum-area measure was developed by Duplantier and Sheffield
\cite{DuplantierSheffieldKPZ2011}. Duplantier, Miller, and Sheffield
identified Brownian contours with
curve-decorated quantum surfaces \cite{DMS}, and Miller and Sheffield
proved the corresponding fixed-area sphere identification
\cite{MillerSheffieldSphere}. These results determine the encoding and
surface limits. The graph metric needs additional control of local
passages and of vertices which typical samples might miss.

The construction of the subcritical LQG metric rests on the LFPP
tightness theorem of Ding, Dub\'edat, Dunlap, and Falconet
\cite{DDDFLFPPTightness}, the weak-metric theory of Dub\'edat, Falconet,
Gwynne, Pfeffer, and Sun \cite{DFGPSWeakMetric}, and the uniqueness
theorem of Gwynne and Miller \cite{GwynneMillerMetric}. Our comparison
argument draws on annular independence and covering arguments for
GFF-local metrics \cite[Lemma~3.1 and Section~4]{GwynneMillerLocalMetrics},
and on the local-field-perturbation and geodesic-saving strategy of the
uniqueness proof. Here the additional task is to obtain and compare
conditional passage laws from discrete maps.

Gwynne and Miller's graph-metric conjecture explicitly includes FK
weights for \(\gamma\in(\sqrt2,2)\) and spanning-tree weights at
\(\gamma=\sqrt2\) \cite[Conjecture~1.7]{GwynneMillerMetric}.
Theorem~\ref{thm:main} resolves its specified finite spherical FK family
and its finite spherical spanning-tree case, with the additional
measured convergence stated above. It does not address every ensemble
or version of that conjecture. For several models with independent-increment
encodings, including infinite spanning-tree-weighted maps, Gwynne,
Holden, and Sun coupled finite contour portions with portions of
mated-CRT maps, obtaining rough isometries with polylogarithmic distortion
\cite[Theorem~1.5]{GwynneHoldenSunGraphDistances}.
Ding and Gwynne's volume-growth estimates
\cite[Theorem~1.6, model~2]{DingGwynneDimension} enter our spanning-tree
proof. For the infinite spanning-tree map, Gwynne and Pfeffer identified
the internal-diameter exponent \cite[Theorem~3.1]{GwynnePfefferDLA};
their volume-counting argument \cite[Lemma~3.2]{GwynnePfefferDLA} is
adapted here to endpoint moments.

At \(q=1\), all \(2^n\) edge subsets of a rooted \(n\)-edge map have
the same weight, so the map marginal is uniform. Following the
Brownian-map scaling-limit and uniqueness theorems of Le Gall and
Miermont \cite{LeGallBrownianMap,MiermontBrownianMap}, its metric limit
is known from the uniform-map theorem
\cite[Corollary~1.2]{BettinelliJacobMiermontUniform} and the identification
of the Brownian sphere with \(\sqrt{8/3}\)-LQG
\cite[Theorem~1.4]{MillerSheffieldBrownianII}
\cite[Corollary~1.4]{GwynneMillerMetric}.
At \(q=2\), Equation~\eqref{eq:q-gamma} gives \(\gamma=\sqrt3\), so
the FK--Ising ensemble is included in the same theorem.
The spanning-tree clause is proved independently and is not obtained
by taking a limit of the FK parameter.

\subsection{Proof overview}

The common objective is to identify local graph distances, control every
vertex, and then retain the measure under the passage to the limit.
The two encodings require different estimates for the first two steps.

\paragraph{Local information from words.}
A contour records a decorated map by a finite word or walk. Buffered
compact regions admit finite chart descriptions containing retained word
intervals, exact boundary heights, and gluing instructions. Refining
these charts recovers every confined passage and its finite marked
subpaths. For FK words, flexible orders can depend on earlier letters
outside these intervals. We retain the finite supplier-comparison
instructions and the order of the relevant traversals, then smooth
relative heights on the exact lattice before passing to the limit.
This proves conditional laws for local passage observations,
including discontinuous shortest-path costs
(Sections~\ref{sec:protected-charts}--\ref{sec:local-kernels}).
The tree proof supplies its own reconstruction and conditional laws from
independent increments in Section~\ref{tree:sec:certificates}.

\paragraph{Identification by local field changes.}
We prove a common comparison and rigidity theorem for local passage
laws. FK maps supply its normalization through surrounding circuits
before arbitrary endpoints are controlled; spanning-tree maps supply
it through internal moment bounds. Annuli at prescribed quantum scales
compare these laws with \(D_h\). If the optimal lower
and upper comparison constants differed, a local change of the real
field would create a strict shortcut near a saturated geodesic segment.
A weighted count of the possible output passages gives a contradiction.
The established Weyl rule controls \(D_h\); exact local conditional laws
transport the graph-passage observations. This establishes a
\emph{deterministic} multiplier at the level of the whole limiting law.
Each model verifies the local-kernel, noncollapse, and joining
hypotheses separately.

\paragraph{Uniform control of endpoints.}
For FK maps, the Brownian two-wall exponent satisfies
\(2/\gamma^2>1/2\). This strict margin pays for the choices of locations
in finite contact diagrams. The resulting entropy estimate and
protected planar sewing produce small surrounding circuits; a recursive
covering argument then controls arbitrary discrete endpoints
(Sections~\ref{arm:section}--\ref{spike:section}).
For spanning trees, internal moment bounds at record scales and the
\(\SLE_8\) interval topology give the endpoint modulus. A first-failure
argument extends these bounds to every scale
(Sections~\ref{tree:sec:moments}--\ref{tree:sec:bootstrap}).

\paragraph{Finite spheres and their area measures.}
The FK local laws transfer to ordinary sphere charts with their actual
field normalization. Two conditionally independent uniform corner roots
of the same map, together with one common sequence of corner samples,
exclude vertices escaping all finite sample nets.
For spanning trees, explicit excursion densities control the two ends
of the contour and give uniform convergence of the metric pulled back
to contour time. The encodings identify the discrete corner measure
with the pushforward of uniform time, and the continuum curve is
parameterized by quantum area. The sampling and common-parameter lemmas
of Section~\ref{measure:section} therefore yield measured convergence
in Sections~\ref{finite:section} and~\ref{tree:sec:sphere}.
A single elementary normalization lemma then removes the deterministic
subsequential multiplier using the unconditional law of log diameter.

Several ingredients are useful independently of the final ensembles:
the transfer of discontinuous local observations by exact lattice
smoothing, comparison by circuits before endpoint tightness, the
balance between contact rank and the entropy of chosen locations, and
the two-root completion argument. The measured convergence criterion
separates finite distance arrays from the additional covering condition
needed to control all points of a compact space.
